% © 2026 Ing. David Jaroš — CC BY-NC-ND 4.0
%
% This work is licensed under a Creative Commons Attribution-NonCommercial-NoDerivatives
% 4.0 International License (CC BY-NC-ND 4.0).
% See LICENSE.md for full license text.

\documentclass[11pt,a4paper]{article}
\usepackage{amsmath,amssymb,amsthm}
\usepackage{mathrsfs}
\usepackage{booktabs}
\usepackage{hyperref}
\usepackage{geometry}
\geometry{margin=2.5cm}

\newtheorem{proposition}{Proposition}
\newtheorem{lemma}{Lemma}
\newtheorem{remark}{Remark}
\newtheorem{gap}{Open Gap}

\newcommand{\proved}[1]{\textbf{[PROVED]} #1}
\newcommand{\heuristic}[1]{\textbf{[HEURISTIC]} #1}
\newcommand{\speculative}[1]{\textbf{[SPECULATIVE]} #1}
\newcommand{\cond}[1]{\textbf{[CONDITIONAL]} #1}
\newcommand{\invalid}[1]{\textbf{[INVALID]} #1}
\newcommand{\plausible}[1]{\textbf{[PLAUSIBLE]} #1}
\newcommand{\derived}[1]{\textbf{[DERIVED]} #1}

\title{Higher-Loop and Threshold Corrections to the $B$-Coefficient}
\author{Ing.\ David Jaro\v{s}\\
\small \texttt{research\_tracks/rg\_B46/}}
\date{May 2026}

\begin{document}
\maketitle

\begin{abstract}
We systematically compute all candidate corrections $\Delta B_i$ that could move the
one-loop baseline $B_0 = 8\pi \approx 25.133$ toward the current KK+winding estimate
$B_\mathrm{KK} \approx 43.6$ and ultimately the phenomenological target $B \approx 46$.
Sources examined: two-loop RG, KK threshold tower, compactification threshold,
winding degeneracy, zero-mode subtraction, ghost determinant, curvature/torsion
corrections, and self-dual torus shape correction.  Each $\Delta B_i$ is classified
as DERIVED, PLAUSIBLE, SPECULATIVE, or INVALID.  \textbf{Hard rule}: no correction
is chosen because it improves the numerical match.
\end{abstract}

\tableofcontents

%------------------------------------------------------------
\section{Setup and Conventions}
%------------------------------------------------------------

The one-loop baseline (from \texttt{one\_loop\_rg\_derivation.tex}) is:
\begin{equation}
  B_0 = \frac{2\pi N_\mathrm{eff}}{3} = 8\pi \approx 25.133,
  \quad N_\mathrm{eff} = 12.
  \label{eq:B0}
\end{equation}
The KK+winding heuristic (from prior work in \texttt{rg\_nlogn/}):
\begin{equation}
  B_\mathrm{KK+wind} \approx \frac{2n}{\pi} \approx 43.6
  \quad (n = 137),
  \label{eq:BKK}
\end{equation}
which is an $n$-dependent formula, heuristically applicable near $n = 137$.
The phenomenological target is $B \approx 46$ (not used as input; compared only at the end).

The total B is:
\begin{equation}
  B_\mathrm{total} = B_0 + \sum_i \Delta B_i.
\end{equation}

\textbf{Hard rule}: $\Delta B_i$ is derived from its stated physical source and
assessed on that basis.  The match (or mismatch) to $B \approx 46$ is stated
only after each $\Delta B_i$ is computed independently.

%------------------------------------------------------------
\section{Two-Loop RG Correction}
%------------------------------------------------------------

\subsection{Derivation}

The two-loop beta function for U(1) with $N_\mathrm{eff} = 12$ complex scalars in 4D is
\begin{equation}
  \beta(e^2)\big|_\mathrm{2-loop}
  = \frac{N_\mathrm{eff}^2}{16\pi^4}\,e^6
  = \frac{144}{16\pi^4}\,e^6
  = \frac{9}{\pi^4}\,e^6.
\end{equation}
The two-loop correction to the running coupling:
\begin{equation}
  \delta e^2_\mathrm{2-loop}(n)
  = -\frac{N_\mathrm{eff}^2}{16\pi^4}\,e^6_0\,\ln n + \mathcal{O}(e^8).
\end{equation}
The induced correction to $B$ (using $B = V(n)/\partial_n^2 V - 1/\partial_n^2 V$, approximately):
\begin{equation}
  \Delta B_\mathrm{2-loop}
  = \frac{N_\mathrm{eff}^2\,e^4_0}{8\pi^4}
  \approx \frac{144 \times g_\psi^4}{8\pi^4},
\end{equation}
where $g_\psi$ is the effective coupling of the $\psi$-circle gauge field.
For $g_\psi \sim 1/\sqrt{4\pi}$ (perturbative coupling):
\begin{equation}
  \Delta B_\mathrm{2-loop}
  \approx \frac{144}{8\pi^4} \times \frac{1}{16\pi^2}
  \approx \frac{1}{88\pi^2}
  \approx 0.0012.
\end{equation}

In the heuristic estimate from \texttt{rg\_nlogn/b\_coefficient\_analysis.md},
$\Delta B_\mathrm{2-loop} \approx 0.05$ (using $g^2 \approx 0.3$).

\subsection{Classification}

\heuristic{$\Delta B_\mathrm{2-loop} \approx 0.05$}: Negligible.

\begin{center}
\begin{tabular}{lccc}
\toprule
Source & Formula & Value & Classification \\\midrule
Two-loop RG & $N_\mathrm{eff}^2 e^4/(8\pi^4)$ & $\lesssim 0.05$ & \textbf{PLAUSIBLE} \\\bottomrule
\end{tabular}
\end{center}

\noindent
\textbf{Conclusion}: Two-loop correction is physically well-defined but
contributes at most $\Delta B \lesssim 0.05$, entirely negligible relative to the
gap $B \approx 46 - 25.1 \approx 20.9$.

%------------------------------------------------------------
\section{KK Threshold Tower}
%------------------------------------------------------------

\subsection{Setup}

The KK tower on $S^1_\psi$ generates a threshold contribution from all modes
with $m_k^2 = k^2$ for $k = 1, 2, \ldots, n-1$ that are integrated out below
the scale $n$.

\subsection{Derivation}

The one-loop sum over all KK modes below scale $n$:
\begin{equation}
  \delta V_\mathrm{KK}(n)
  = -\frac{N_\mathrm{eff}}{4\pi}\sum_{k=1}^{n-1} k^2
    \left(\ln\frac{k^2}{\mu^2} - 1\right).
\end{equation}
Using the Euler–Maclaurin formula for large $n$:
\begin{equation}
  \sum_{k=1}^{n-1} k^2 \ln k
  \approx \int_1^n dk\, k^2 \ln k
  = \frac{n^3\ln n}{3} - \frac{n^3}{9} + \mathcal{O}(n^2\ln n).
\end{equation}
Therefore:
\begin{equation}
  \delta V_\mathrm{KK}(n)
  \approx -\frac{N_\mathrm{eff}}{4\pi}\cdot\frac{n^3\ln n}{3}
  = -\frac{N_\mathrm{eff}\,n^3\ln n}{12\pi}.
\end{equation}
This is an $n^3\ln n$ correction, \textbf{not} $n\ln n$.
It renormalises higher-order terms but does not directly shift $B$.

However, the \emph{threshold correction} from integrating out the mode at scale $k$
one at a time contributes a step-function correction to the running coupling:
\begin{equation}
  \delta e^2(n) = -\frac{N_\mathrm{eff}}{12\pi^2}\,\sum_{k=1}^{n}\theta(n-k),
\end{equation}
which counts $n$ threshold crossings.  This gives an additional contribution:
\begin{equation}
  \Delta B_\mathrm{KK-thresh}
  = \frac{N_\mathrm{eff}\,n}{12\pi^2}
  = \frac{12n}{12\pi^2}
  = \frac{n}{\pi^2}.
\end{equation}
At $n = 137$: $\Delta B_\mathrm{KK-thresh} = 137/\pi^2 \approx 13.9$.

\heuristic{$\Delta B_\mathrm{KK-thresh} \approx 13.9$}: Physical (step-function matching)
but depends on $n = 137$ being fixed, making $B$ $n$-dependent.

\begin{center}
\begin{tabular}{lccc}
\toprule
Source & Formula & Value & Classification \\\midrule
KK threshold tower & $N_\mathrm{eff}\,n/(12\pi^2)$ & $\approx 13.9$ & \textbf{PLAUSIBLE} \\\bottomrule
\end{tabular}
\end{center}

%------------------------------------------------------------
\section{Compactification Threshold}
%------------------------------------------------------------

\subsection{Derivation}

At the compactification scale $M_\mathrm{comp} = 1/R_\psi = 1$ (natural units),
the 5D theory matches to the 4D effective theory.  The threshold correction is:
\begin{equation}
  \Delta e^2\big|_\mathrm{comp}
  = -\frac{N_\mathrm{eff}}{12\pi^2}\ln\frac{M_\mathrm{comp}^2}{\mu^2}
  = -\frac{N_\mathrm{eff}}{12\pi^2}\ln\frac{1}{\mu^2}
  = \frac{N_\mathrm{eff}}{6\pi^2}\ln\mu.
\end{equation}
At $\mu = 1$ this vanishes.

\heuristic{The compactification threshold is scheme-dependent and vanishes in
$\overline{\mathrm{MS}}$ at $\mu = M_\mathrm{comp}$.}

\begin{equation}
  \Delta B_\mathrm{comp} = 0 \quad (\overline{\mathrm{MS}},\; \mu = M_\mathrm{comp}).
\end{equation}

\begin{center}
\begin{tabular}{lccc}
\toprule
Source & Formula & Value & Classification \\\midrule
Compactification threshold & 0 in $\overline{\mathrm{MS}}$ at $\mu=M_\mathrm{comp}$ & $0$ & \textbf{DERIVED} \\\bottomrule
\end{tabular}
\end{center}

%------------------------------------------------------------
\section{Winding Degeneracy (T-Duality Factor)}
%------------------------------------------------------------

\subsection{Derivation}

At the self-dual radius $R_\psi = 1$, T-duality maps KK modes to winding modes.
Both contribute equally to the one-loop vacuum polarisation.  The total (KK +
winding) contribution is:
\begin{equation}
  B_\mathrm{KK+wind}
  = B_\mathrm{KK,one-loop} \times 2
  = \frac{N_\mathrm{eff}\,n}{6\pi^2} \times 2
  = \frac{N_\mathrm{eff}\,n}{3\pi^2}.
\end{equation}
At $n = 137$, $N_\mathrm{eff} = 12$:
\begin{equation}
  B_\mathrm{KK+wind} = \frac{12 \times 137}{3\pi^2}
  = \frac{1644}{3\pi^2}
  \approx \frac{1644}{29.61}
  \approx 55.5.
\end{equation}

\begin{remark}
The existing estimate $B_\mathrm{KK+wind} \approx 43.6 = 2n/\pi$ uses a \emph{different}
one-loop formula ($B_\mathrm{1-loop} = n/(2\pi)$, from the 1D effective theory).
The formula $N_\mathrm{eff}\,n/(3\pi^2)$ from the 4D one-loop approach gives
$\approx 55.5$ at $n=137$.  Neither matches $B \approx 46$ exactly; both are
\textbf{$n$-dependent heuristics}.
\end{remark}

\heuristic{$B_\mathrm{KK+wind}$:
\begin{itemize}
  \item Self-dual radius $R_\psi = 1$ is not derived from UBT moduli; assumed.
  \item T-duality at $R_\psi = 1$ doubles the mode count; standard string result.
  \item The doubling factor 2 is valid if winding modes have the same one-loop
        coefficient as KK modes (symmetric spectrum). This is T-duality assumption.
\end{itemize}}

\begin{center}
\begin{tabular}{lccc}
\toprule
Source & Formula & Value at $n=137$ & Classification \\\midrule
Winding degeneracy (1D baseline) & $2 \times n/(2\pi)$ & $43.6$ & \textbf{HEURISTIC} \\
Winding degeneracy (4D baseline) & $2 \times N_\mathrm{eff}\,n/(6\pi^2)$ & $55.5$ & \textbf{PLAUSIBLE} \\\bottomrule
\end{tabular}
\end{center}

The increment from $B_0 = 8\pi \approx 25.1$:
\begin{equation}
  \Delta B_\mathrm{wind} = B_\mathrm{KK+wind} - B_0 \approx 43.6 - 25.1 = 18.5
  \quad\text{(1D baseline)}.
\end{equation}

%------------------------------------------------------------
\section{Zero-Mode Subtraction}
%------------------------------------------------------------

\subsection{Derivation}

The $n=0$ KK mode is the $\psi$-independent field $\Theta_0(x)$.  It is neutral
under $\mathrm{U}(1)_\psi$ and does not contribute to $\Pi_{\mu\nu}$.

In the sum over KK modes, the zero mode can produce a spurious contribution if the
summation range is extended to include $k = 0$.  A careful treatment excludes it:
\begin{equation}
  \Delta B_\mathrm{zero-mode} = 0 \quad \text{(correctly excluded from $\Pi$)}.
\end{equation}

\derived{$\Delta B_\mathrm{zero-mode} = 0$}: The zero mode is neutral; its exclusion
from $\Pi$ is exact.

\begin{center}
\begin{tabular}{lccc}
\toprule
Source & Formula & Value & Classification \\\midrule
Zero-mode subtraction & 0 (neutral) & $0$ & \textbf{DERIVED} \\\bottomrule
\end{tabular}
\end{center}

%------------------------------------------------------------
\section{Ghost Determinant}
%------------------------------------------------------------

\subsection{Derivation}

For the $\mathrm{U}(1)_\psi$ gauge field, Faddeev--Popov ghosts are non-interacting
complex scalars with fermionic statistics and charge 0.  Their contribution to the
vacuum polarisation:
\begin{equation}
  \Pi^\mathrm{ghost}(k^2) = 0 \quad \text{(ghosts uncharged under U(1)$_\psi$)}.
\end{equation}

For the non-abelian sectors embedded in $\mathcal{B}$:
ghosts subtract the longitudinal polarisation of gauge bosons.  Since the
$N_\mathrm{eff} = 12$ mode count already counts only transverse modes (§2.7 of
the N\_eff audit), ghost subtraction is already accounted for.

\derived{$\Delta B_\mathrm{ghost} = 0$ for U(1)$_\psi$; ghost subtraction is
already incorporated in $N_\mathrm{eff} = 12$ (transverse modes only).}

\begin{center}
\begin{tabular}{lccc}
\toprule
Source & Formula & Value & Classification \\\midrule
Ghost determinant & 0 (already subtracted) & $0$ & \textbf{DERIVED} \\\bottomrule
\end{tabular}
\end{center}

%------------------------------------------------------------
\section{Curvature and Torsion Corrections}
%------------------------------------------------------------

\subsection{Derivation}

If the background metric $g_{\mu\nu}$ has non-zero curvature $R \neq 0$, the
minimally coupled one-loop effective action receives a conformal non-minimal coupling:
\begin{equation}
  S \supset -\xi R \int d^4x\,|\Theta|^2,
  \quad \xi_\mathrm{minimal} = 0,\quad \xi_\mathrm{conformal} = \frac{1}{6}.
\end{equation}
The effective mass is shifted: $m_n^2 \to n^2 + \xi R$.  For $\xi = 0$
(minimal coupling, consistent with UBT default):
\begin{equation}
  \Delta B_\mathrm{curv} = 0 \quad (\xi = 0,\; R = 0\;\text{background}).
\end{equation}
For $\xi = 1/6$ (conformal coupling):
\begin{equation}
  \Delta B_\mathrm{curv}(\xi=1/6)
  = \frac{N_\mathrm{eff}}{6}\cdot\frac{R}{n}\cdot\frac{1}{4\pi^2}.
\end{equation}
For a flat background $R = 0$: $\Delta B_\mathrm{curv} = 0$.

If torsion $T_{\mu\nu\rho} \neq 0$ is present (as possible in UBT with biquaternion
connection), the fermion sector couples to torsion via the spin connection.  For the
bosonic field $\Theta$ in the absence of intrinsic spin coupling:
\begin{equation}
  \Delta B_\mathrm{torsion} = 0 \quad \text{(bosonic scalar, torsion decouples)}.
\end{equation}

\speculative{In a full UBT background with $R \neq 0$ and $T \neq 0$, these
contributions are non-zero.  Their magnitude depends on the specific background
geometry, which is not specified here.}

\begin{center}
\begin{tabular}{lccc}
\toprule
Source & Formula & Value (flat) & Classification \\\midrule
Curvature ($\xi=0$) & $0$ & $0$ & \textbf{DERIVED} \\
Curvature ($\xi=1/6$) & $N_\mathrm{eff} R/(24\pi^2 n)$ & $0$ (flat) & \textbf{SPECULATIVE} \\
Torsion (bosonic $\Theta$) & $0$ & $0$ & \textbf{DERIVED} \\\bottomrule
\end{tabular}
\end{center}

%------------------------------------------------------------
\section{Self-Dual Torus Shape Correction}
%------------------------------------------------------------

\subsection{Derivation}

The compact $\psi$-direction is modelled as $S^1_{R_\psi}$.  A shape correction
arises if the torus has a modular parameter $\tau_\psi = i R_\psi + \epsilon$
with $\epsilon \neq 0$ (off-diagonal component; complex structure).

For the rectangular torus ($\epsilon = 0$, $R_\psi = 1$): no shape correction.

For a self-dual torus at $\tau_\psi = i$ (purely imaginary, no off-diagonal):
the spectrum is symmetric under $n \to -n$ and the shape correction is absent.

If $\epsilon \neq 0$, the eigenvalues of $-\partial_\psi^2$ are modified:
\begin{equation}
  m_{n,m}^2 = \frac{|n + m\tau_\psi|^2}{\mathrm{Im}\,\tau_\psi^2},
\end{equation}
which breaks the $n \to -n$ symmetry and introduces corrections of order
$|\epsilon|^2 n^2$ to the spectrum.  These are:
\begin{equation}
  \Delta B_\mathrm{torus} = \mathcal{O}(|\epsilon|^2).
\end{equation}
For $\tau_\psi = i$ exactly (assumed): $\Delta B_\mathrm{torus} = 0$.

\speculative{If the UBT moduli stabilise the torus at $\tau_\psi \neq i$, a
non-zero shape correction would arise.  Its derivation requires solving the
moduli problem, which is open.}

\begin{center}
\begin{tabular}{lccc}
\toprule
Source & Formula & Value & Classification \\\midrule
Self-dual torus shape ($\epsilon = 0$) & $0$ & $0$ & \textbf{CONDITIONAL} \\
Self-dual torus shape ($\epsilon \neq 0$) & $\mathcal{O}(|\epsilon|^2)$ & unknown & \textbf{SPECULATIVE} \\\bottomrule
\end{tabular}
\end{center}

%------------------------------------------------------------
\section{Summary Table of All Corrections}
%------------------------------------------------------------

\begin{center}
\begin{tabular}{llccc}
\toprule
Source & $\Delta B_i$ & Value & Classification & Included? \\\midrule
Two-loop RG & $N_\mathrm{eff}^2 e^4/(8\pi^4)$ & $\lesssim 0.05$ & PLAUSIBLE & \checkmark \\
KK threshold tower & $N_\mathrm{eff}\,n/(12\pi^2)$ & $\approx 13.9$ & PLAUSIBLE & \checkmark \\
Compactification threshold & 0 (in $\overline{\mathrm{MS}}$) & 0 & DERIVED & \checkmark \\
Winding degeneracy (T-dual) & $\approx B_0$ (adds factor 2) & $\approx 18.5$ & HEURISTIC & \checkmark \\
Zero-mode subtraction & 0 & 0 & DERIVED & \checkmark \\
Ghost determinant & 0 & 0 & DERIVED & \checkmark \\
Curvature ($\xi=0$, flat) & 0 & 0 & DERIVED & \checkmark \\
Torsion (bosonic) & 0 & 0 & DERIVED & \checkmark \\
Self-dual torus ($\epsilon=0$) & 0 & 0 & CONDITIONAL & \checkmark \\\bottomrule
\end{tabular}
\end{center}

\noindent
Dominant non-zero corrections (dependent on $n$, evaluated at $n \approx 137$):
\begin{itemize}
  \item Winding degeneracy: $\Delta B_\mathrm{wind} \approx 18.5$ (heuristic; T-duality at self-dual radius).
  \item KK threshold tower: $\Delta B_\mathrm{KK} \approx 13.9$ (plausible; step-function matching).
\end{itemize}

\noindent
The two dominant corrections are \textbf{not independent}: the winding
degeneracy already accounts for most of the KK threshold contribution.
Using \emph{either one} (not both) gives:
\[
  B_\mathrm{total}^\mathrm{wind}
  = B_0 + \Delta B_\mathrm{wind}
  = 8\pi + 18.5
  \approx 43.6,
\]
which recovers the prior KK+winding estimate.  Adding both overcounts.

\begin{gap}[Remaining discrepancy $\Delta B \approx 2.4$]
\label{gap:2.4}
After the winding correction, $B_\mathrm{total} \approx 43.6$.
The residual gap $\Delta B \approx 2.4$ to reach $B \approx 46$
is \textbf{not accounted for} by any of the computed corrections.
Two-loop: $\lesssim 0.05$; curvature: 0 (flat); ghost: 0; zero-mode: 0; torus shape: 0.
The gap $\Delta B \approx 2.4$ ($\approx 5.2\%$ of $B$) is unresolved.
\end{gap}

\end{document}
